\chapter*{Agradecimientos}
\addcontentsline{toc}{chapter}{Agradecimientos}

A la Universidad Nacional de Ingeniería, casa de estudios que forjó mi formación
profesional y me brindó las herramientas intelectuales para enfrentar los desafíos
de la ingeniería moderna. Su tradición de excelencia académica y su compromiso
con el desarrollo científico y tecnológico del país son una inspiración permanente.

A mis profesores de la Facultad de Ingeniería Industrial y de Sistemas, quienes con
su dedicación, rigor académico y generosidad compartieron no solo conocimientos,
sino también una forma de pensar y resolver problemas. En especial, mi profundo
agradecimiento a mi asesor de tesis, por su orientación, paciencia y valiosas
observaciones que enriquecieron cada etapa de esta investigación.

A mis hijos, Harvey y Degel, que son mi mayor motivación y el recordatorio más
poderoso de que el esfuerzo de hoy construye el país de mañana. En ustedes veo
el futuro del Perú: su curiosidad, su energía y su enorme potencial me impulsan a
ser mejor profesional y mejor persona cada día. Esta tesis es, en gran parte, para
ustedes.

A todas las personas del sector bancario que participaron en el proceso Delphi y
contribuyeron con su experiencia y criterio profesional al desarrollo de esta
metodología. Su tiempo y conocimiento fueron fundamentales para dotar de rigor y
validez a este trabajo.
